\chapter{界面与日常编辑}\label{ch:workspace}
\section{常用入口在哪里}
\par\noindent\begin{tabularx}{\linewidth}{@{}P{38mm}Y@{}}
\toprule
入口 & 用途\\\midrule
\ui{文件}菜单 & 打开项目、导入 Markdown/文本、保存源码、打开恢复副本。\\
\ui{文件 → 保存项目} & 保存原生项目；第一次保存时选择位置。\\
\ui{编辑 → 撤销 / 重做} & 撤销或恢复正文编辑。评价有另外的撤销按钮。\\
顶部\ui{设置} & 打开右侧文档设置，分为正文、元素、字形、问题；保留上次选择的分类。\\
顶部\ui{字形评价} & 展开或收起底部评价区。\\
顶部\ui{导出 SVG…} & 输出完整文档的写字机文件。旁边下拉菜单提供参考页面与评价 CSV 导出。\\
底部\ui{查看问题} & 查看当前排版诊断，定位错误或模板问题。\\
\bottomrule
\end{tabularx}

左侧源码、右侧预览是主要工作区。拖动中间分隔线可以分配宽度；不需要的设置或评价面板可以收起。关闭设置面板不会取消已经确认的文档设置。

\section{编辑、定位与撤销}
直接输入或粘贴源码后，预览稍候更新。输入法的候选和未提交拼音不应成为正文。重排时可能暂时保留旧页，等到“预览已更新”后再判断最终结果。

点击预览中的字可定位相应源码并选中具体字形。点击任务方框会改变源码中的勾选状态，支持撤销。图片点击定位整段图片源码，图片本身不作为字形评分。

\par\noindent\begin{tabularx}{\linewidth}{@{}P{31mm}P{40mm}Y@{}}
\toprule
动作 & macOS & Windows\\\midrule
打开项目 & Command+O & Ctrl+O\\
保存项目 & Command+S & Ctrl+S\\
剪切 / 复制 / 粘贴 & Command+X/C/V & Ctrl+X/C/V\\
全选 & Command+A & Ctrl+A\\
撤销 & Command+Z & Ctrl+Z\\
重做 & Command+Shift+Z & 以\ui{编辑}菜单为准\\
\bottomrule
\end{tabularx}

\ui{编辑}菜单的操作发给当前有焦点的输入框。若正在编辑设置中的数字或文字，撤销与粘贴也可能作用于该控件；需要操作正文时，先点击源码区。源码标题旁的\ui{粘贴}按钮会明确把剪贴板文本插入正文光标处。

\section{正文设置}
\entry{设置 → 正文}
\par\noindent\begin{tabularx}{\linewidth}{@{}P{33mm}P{23mm}Y@{}}
\toprule
参数 & 新文档默认 & 调整的效果\\\midrule
字高 & 6 mm & 控制普通文字的标称字高；实际墨迹高度由字形决定。\\
行高 & 9 mm & 控制普通行的基线间距；复杂内容可能需要更高的行。\\
笔宽 & 0.3 mm & 控制统一描线宽度，不模拟原笔压变化。\\
段距 / 标题段距 & 2 / 3 mm & 调节块与块之间的空隙。\\
首行缩进 & 0 mm & 控制普通段落的首行起点。\\
列表 / 引用缩进 & 8 / 4 mm & 调节相应内容的缩进。\\
\bottomrule
\end{tabularx}

通过\ui{标题字号…}设置各级标题的倍率。字号变大可能改变断行和分页，重新检查表格、公式和模板区域。不要用缩放预览代替调整实际字高。

\section{选择内容模式}
\ui{Markdown 混排}解释标题、表格和公式等语法；\ui{纯文本检查}用于按字面文字检查字形。Markdown 格式选择只在混排下有意义。已有项目打开时会恢复自己的模式与设置，不一定与新建默认值相同。
